% ============================================================================
%  অধ্যায় ১ — শুরুর আগে
%  COMPILE WITH XeLaTeX  (twice, so the point total resolves)
%  Overleaf: Menu -> Compiler -> XeLaTeX
% ============================================================================
\documentclass[11pt]{article}
\usepackage{bnhandout}          % add [nosol] for a student edition

\hauthor[https://jugantordey.github.io/]{Jugantor Dey}
\hseries{\textit{Mathematics for the Physics Olympiad}}

\begin{document}

\handouttitle{Chapter 1 : Before You Dive In}

\noindent
এই অধ্যায়ের কোনো পূর্বশর্ত নেই। অধ্যায় ০ আর এইটা পরস্পর স্বাধীন :
ওখানে সংখ্যা নিয়ে কথা, এখানে যুক্তি নিয়ে। যেকোনো ক্রমে পড়া যায়।

\noindent
এই অধ্যায়ে একটাও নতুন সূত্র নেই, একটাও হিসাব নেই। তবু এটাই সিরিজের সবচেয়ে
গুরুত্বপূর্ণ অধ্যায়, কারণ এখানে ঠিক করা হচ্ছে বাকি সবগুলো কীভাবে পড়তে হবে।

\noindent
স্কুলে গণিত মানে সাধারণত : নিয়ম মুখস্থ করো, নিয়ম প্রয়োগ করো, উত্তর মেলাও।
এই সিরিজে সেটা চলবে না। এখানে প্রতিটি সূত্র কোথা থেকে এলো তা দেখানো হবে, আর
তোমার কাছেও চাওয়া হবে ``কেন'' -এর উত্তর, শুধু ``কত'' -এর উত্তর নয়। সেই বদলটার জন্য
তিনটে জিনিস আগে বুঝে নেওয়া দরকার : প্রমাণ জিনিসটা আসলে কী, গণিতের চিহ্নগুলো
সত্যিকারে কী বোঝায়, আর কোনো যুক্তি কখন নিখুঁত আর কখন ফাঁকি। এখানে মোট
\totalpoints\ পয়েন্ট আছে।

\section{কেন গণিত শিখব}

\noindent
পদার্থবিজ্ঞানের জন্য গণিত লাগে, এ কথাটা সবাই বলে। কিন্তু কেন লাগে, সেটা
সাধারণত ভুলভাবে বলা হয়। বলা হয়, হিসাব করার জন্য লাগে। আসল কারণটা অন্য।

\begin{idea}[গণিত হিসাবের যন্ত্র নয়, নিখুঁত করে বলার ভাষা]
বাংলায় বা ইংরেজিতে একটা কথা বললে সেটা প্রায়ই একাধিক অর্থ বহন করে। ``ভারী
জিনিস দ্রুত পড়ে'' কথাটা ধরো। কতটা ভারী? কতটা দ্রুত? বাতাস আছে কি নেই? কথাটা
সত্য না মিথ্যা, তা যাচাই করার আগেই বোঝা দরকার কথাটা ঠিক কী দাবি করছে; আর
এই বাক্যটা তা বলছে না।

এখন লেখো : বাতাসশূন্য ঘরে স্থির অবস্থা থেকে ছাড়া যেকোনো বস্তুর $2$ সেকেন্ড সময় পরে
পতনের দূরত্ব $19.6$ মিটার,যা বস্তুর ভরের উপর নির্ভর করে না।

এবার আর দ্বিতীয় কোনো অর্থ নেই। কথাটা এখন হয় সত্য, নয় মিথ্যা, আর পরীক্ষা করে
দেখা যায় কোনটা। গণিতের আসল কাজ এটাই : দাবিকে এমন আকারে আনা যাতে তা ঠিক বা ভুল প্রমাণ করা সম্ভব হয়।
\end{idea}

\begin{remark}
যে দাবি কোনো অবস্থাতেই ভুল হতে পারে না, সে দাবি কোনো তথ্যও দেয় না। ``আগামীকাল
বৃষ্টি হতেও পারে, নাও হতে পারে'' কথাটা কখনো মিথ্যা হবে না, আর ঠিক সেজন্যই কথাটা
অর্থহীন। পদার্থবিজ্ঞানে একটা তত্ত্বের মূল্য মাপা হয় সে কত নির্দিষ্ট করে বলছে
তা দিয়ে, কারণ যত নির্দিষ্ট, তত সহজে ধরা পড়বে যদি সে ভুল হয়।
\end{remark}

\section{প্রমাণ কী}

\noindent
তিনটে শব্দ প্রায়ই গুলিয়ে ফেলা হয়। আলাদা করে নিই।

\begin{idea}[বিবৃতি, দাবি, প্রমাণ]
\textbf{বিবৃতি (statement)} : এমন একটি বাক্য যা হয় সত্য নয় মিথ্যা।
``$7$ একটি মৌলিক সংখ্যা'' একটি বিবৃতি। ``$x + 3$'' বিবৃতি নয়, কারণ ওটা কোনো তথ্য দিচ্ছে না। ``অঙ্কটা সুন্দর'' -ও বিবৃতি নয়, কারণ সত্য-মিথ্যা যাচাইয়ের
কোনো উপায় নেই।

\textbf{দাবি (claim)} : যে বিবৃতিটাকে আমরা সত্য বলে দেখাতে চাইছি, কিন্তু
এখনো দেখাইনি।

\textbf{প্রমাণ (proof)} : ইতিমধ্যে মেনে নেওয়া সত্য থেকে শুরু করে ধাপে ধাপে
দাবিটিতে পৌঁছানোর একটা পথ, যেখানে প্রতিটি ধাপ আগের ধাপ থেকে অনিবার্যভাবে আসে।

মূল কথা : প্রমাণ কোনো একটা উদাহরণ নয়, আর জোরালো বিশ্বাসও নয়। প্রমাণ একটা
\emph{পথ}, যা যে কেউ ধরে হাঁটতে পারে এবং হাঁটলে একই জায়গায় পৌঁছাবে।
\end{idea}

\begin{example}
দেখাও যে দুটি জোড় সংখ্যার যোগফল জোড়।
\end{example}

\begin{solbox}
প্রথমে ``জোড়'' কথাটার একটা কাজে লাগানোর মতো রূপ দরকার। $n$ (ধরে নাও $2, 4, 6, 8$-এর মতোই $n$ একটা সংখ্যা) জোড় মানে $n = 2k$,
যেখানে $k$ একটি পূর্ণসংখ্যা (কারণ যেকোনো জোড় সংখ্যাই 2-এর গুণিতক)।

ধরি $m$ ও $n$ দুটি জোড় সংখ্যা। তাহলে এমন পূর্ণসংখ্যা $j$ ও $k$ আছে যে
$m = 2j$ এবং $n = 2k$। যোগ করি :
\[
  m + n = 2j + 2k = 2(j+k) .
\]
$j + k$ দুটি পূর্ণসংখ্যার যোগফল, তাই নিজেও পূর্ণসংখ্যা। অর্থাৎ $m+n$-কে
$2 \times (\text{পূর্ণসংখ্যা})$ আকারে লেখা গেল, তাই $m+n$ জোড়।

লক্ষ করো, কোথাও কোনো নির্দিষ্ট সংখ্যা বসানো হয়নি। যুক্তিটা \emph{সব} জোড়
সংখ্যার জোড়ার জন্য একসাথে খেটেছে। এটাই সরাসরি প্রমাণ (direct proof)।
\end{solbox}

\noindent
এবার সবচেয়ে দরকারি কথাটায় আসি। প্রমাণ করা আর ভুল প্রমাণ করা, এই দুটো কাজ
একই নয়, একেবারেই একই নয়।

\begin{idea}[উদাহরণ দিয়ে দেখানো প্রমাণ নয়]
একটা দাবি ``সব ক্ষেত্রে'' সত্য বলতে চাইলে কয়েকটা ক্ষেত্রে পরীক্ষা করা কখনোই
যথেষ্ট নয়, তুমি যতগুলো ক্ষেত্রেই পরীক্ষা করো না কেন। কারণ যেগুলো পরীক্ষা করোনি
তাদের সংখ্যা সাধারণত অসীম, আর দাবি ভাঙার জন্য তাদের একটাই যথেষ্ট।

কথাটা কতটা মারাত্মক, একটা উদাহরণে বোঝা যাক। ধরো দাবি :
\[
  n^2 + n + 41 \text{ সবসময় একটি মৌলিক সংখ্যা।}
\]
$n = 0$ বসাও, পাও $41$, মৌলিক। $n = 1$, পাও $43$, মৌলিক। $n = 2$, পাও $47$,
মৌলিক। এভাবে চালিয়ে যাও। $n = 39$ পর্যন্ত \emph{প্রতিটি} মান মৌলিক আসে।
চল্লিশটা সফল পরীক্ষা। যেকোনো পরীক্ষাগারে এটা যথেষ্ট প্রমাণ বলে গণ্য হতো।

তারপর $n = 40$ বসাও :
\[
  40^2 + 40 + 41 = 1600 + 40 + 41 = 1681 = 41 \times 41 .
\]
মৌলিক নয়। দাবিটা মিথ্যা, এবং চল্লিশটা সফল পরীক্ষা তার কোনো রক্ষা করতে পারল না।
\end{idea}

\begin{idea}[একটিমাত্র counterexample একটি দাবিকে দুমড়ে মুচড়ে দিতে পারে!]
উল্টো দিকটা আশ্চর্যরকম সহজ। কোনো দাবি \emph{মিথ্যা} দেখাতে হলে একটামাত্র
ব্যতিক্রম দেখালেই চলে। এই ব্যতিক্রমটার নাম counterexample।

``সব মৌলিক সংখ্যা বিজোড়'' দাবিটা ভাঙতে $2$ দেখানোই যথেষ্ট। আর কিছু বলার দরকার
নেই, কোনো সাধারণ যুক্তি লাগবে না। একটা উদাহরণেই কাজ শেষ।

অর্থাৎ কাজ দুটো অসম :
\begin{enumerate}
  \item দাবি সত্য দেখাতে : সব ক্ষেত্র একসাথে সামলাতে হবে, তাই সাধারণ যুক্তি লাগবে।
  \item দাবি মিথ্যা দেখাতে : একটিমাত্র ক্ষেত্রই যথেষ্ট।
\end{enumerate}
এই অসমতাটা এই সিরিজে বারবার ফিরে আসবে, বিশেষ করে অধ্যায় ২-এ অভেদ ও
সমীকরণের পার্থক্য বোঝার সময়।
\end{idea}

\begin{example}
``$n$ পূর্ণসংখ্যা হলে $n^2 \geq n$'' \; দাবিটি কি সত্য?
\end{example}

\begin{solbox}
$n = 0, 1, 2, 3, \dots$ সবই খাটে, আর ঋণাত্মকগুলোতেও খাটে, কারণ বর্গ অঋণাত্মক
হয়ে যায়। মনে হচ্ছে সত্য, এবং পূর্ণসংখ্যার জন্য সত্যিই সত্য।

কিন্তু শর্তটা থেকে ``পূর্ণসংখ্যা'' শব্দটা তুলে দিলেই দাবিটা ভেঙে পড়ে :
$n = \tfrac{1}{2}$ নিলে $n^2 = \tfrac{1}{4}$, যা $\tfrac{1}{2}$-এর চেয়ে ছোট।

শিক্ষণীয় কথা : দাবির শর্তগুলো সাজসজ্জা নয়। ``পূর্ণসংখ্যা হলে'' অংশটা বাদ দিলে
দাবিটা আর সত্য থাকে না। কোনো উপপাদ্য পড়ার সময় শর্তগুলো আগে পড়ো, কারণ
ওখানেই বলা আছে উপপাদ্যটা কোথায় খাটে আর কোথায় খাটে না।
\end{solbox}

\section{গাণিতিক ভাষা}

\noindent
গণিতের চিহ্নগুলো লেখা ছোট করার জন্য নয়। প্রতিটির একটা নির্দিষ্ট মানে আছে, আর
ভুল চিহ্ন বসালে লেখাটা ভুল কথা বলে। সবচেয়ে বেশি ভুল হয় প্রথম দুটোতে।

\begin{idea}[$\Rightarrow$ ও $\Leftrightarrow$]
$P \Rightarrow Q$ পড়া হয় ``$Q$ সত্য হবে যদি $P$  সত্য হয়''। এর মানে : $P$ সত্য হলে $Q$-ও সত্য
হতে বাধ্য। এটা \emph{উল্টো দিকের} কিছু বলে না।

$P \Leftrightarrow Q$ পড়া হয় ``$Q$ সত্য হবে যদি ও কেবল যদি $P$  সত্য হয়''। এর মানে দুই দিকই :
$P \Rightarrow Q$ এবং $Q \Rightarrow P$।

পার্থক্যটা দেখো :
\[
  x = 2 \;\Rightarrow\; x^2 = 4 .
\]
এটা সত্য। কিন্তু উল্টোটা সত্য নয়, কারণ $x^2 = 4$ হলে $x$ হতে পারে $2$ অথবা
$-2$। তাই এখানে $\Leftrightarrow$ লেখা \emph{ভুল} হতো।

অন্যদিকে
\[
  x - 3 = 0 \;\Leftrightarrow\; x = 3
\]
এখানে দুই দিকই খাটে, তাই $\Leftrightarrow$ সঠিক।
\end{idea}

\begin{remark}
সমীকরণ সমাধানের সময় এই পার্থক্যটাই সবচেয়ে গুরুত্বপূর্ণ। তুমি যদি ধাপে ধাপে
$\Rightarrow$ দিয়ে এগোও, তাহলে শেষে যা পাবে তা ``সম্ভাব্য উত্তর'', নিশ্চিত
উত্তর নয়; তাই যাচাই করা বাধ্যতামূলক। আর প্রতিটি ধাপ যদি $\Leftrightarrow$ হয়,
তবেই কেবল যাচাই না করে ছেড়ে দেওয়া যায়।

দুই পাশে বর্গ করা ধাপটা সবচেয়ে বিখ্যাত ফাঁদ, কারণ ওটা $\Rightarrow$, কখনোই
$\Leftrightarrow$ নয়। এই কারণেই বর্গ করে সমাধান করলে মাঝেমধ্যে এমন ``উত্তর''
বেরোয় যা আসল সমীকরণে বসালে মেলে না।
\end{remark}

\begin{idea}[বাকি চিহ্নগুলো]
\begin{enumerate}
  \item $\therefore$ \; ``সুতরাং''। ফলাফল দেখানোর আগে বসে।
  \item $\because$ \; ``যেহেতু''। কারণ দেখানোর আগে বসে।
  \item $\forall$ \; ``সকল''। $\forall x$ মানে ``সকল $x$-এর জন্য''।
  \item $\exists$ \; ``বিদ্যমান''। $\exists \, x$ মানে ``এমন অন্তত একটি $x$ আছে''।
\end{enumerate}

\noindent
শেষ দুটোর ক্রম বদলালে অর্থ পুরো বদলে যায়, আর এই ভুলটা খুব সাধারণ। ``প্রতিটি
তালার জন্য এমন একটা চাবি আছে যা ওই তালাটা খোলে'' আর ``এমন একটা চাবি আছে যা প্রতিটি
তালা খোলে'' দুটো একেবারে আলাদা দাবি। প্রথমটা সাধারণ ব্যাপার, দ্বিতীয়টা
master-key!
\end{idea}

\begin{remark}
এই সিরিজে $\Rightarrow$ ও $\Leftrightarrow$ প্রায় প্রতিটি অধ্যায়ে লাগবে।
$\therefore\,$, $\because\,$, $\forall$, $\exists$ এখানে খুব একটা ব্যবহার করা হবে
না, কারণ বাংলায় কথাগুলো লিখে দিলেই পড়তে সহজ হয়। তবু চিহ্নগুলো চিনে রাখা
দরকার, কারণ অন্য ম্যাথ বইয়ে, বিশেষ করে ইংরেজি ভাষার বইয়ে, এগুলো অহরহ আসবে।
\end{remark}

\section{কীভাবে এই সিরিজ পড়বে}

\noindent
প্রতিটি অধ্যায়ে কয়েক ধরনের বাক্স আছে, আর প্রতিটির আলাদা কাজ।

\begin{idea}[বাক্সগুলো অর্থ]
\begin{enumerate}
  \item \textbf{ধারণা} : মূল কথা। এখানে যা আছে তা বুঝতেই হবে; মুখস্থ নয়, বুঝতে হবে।
        কোনো ধারণার বাক্স আটকে গেলে সামনে এগোনোর কোনো মানে নেই; আবার পিছন থেকে পড়ে আসতে হবে, বড় কারও সাহায্য নিতে হবে, আমাকে জানাতে হবে (এই অধ্যায় যেই ওয়েবসাইট  থেকে ডাউনলোড করেছ সেখানে লিখে দিয়েছি কিভাবে আমাকে জানাবে), আর একান্তই না পারলে ইন্টারনেট থেকে খুঁজে শিখে নিবে (বিভিন্ন এড-টেক প্ল্যাটফর্মের ভিডিও এড়িয়ে চলাই ভালো, ইন্টারনেট থেকে বই, আর্টিকেল খুঁজে পড়বে, অথবা ইউটিউবে কোনো ইংলিশ চ্যানেল থেকে শিখতে পারো, যেমন professor dave explains বা 3blue1brown ইত্যাদি)।
  \item \textbf{উদাহরণ} ও \textbf{সমাধান} : ধারণাটা কাজে লাগিয়ে দেখানো হচ্ছে।
        সমাধান পড়ার আগে নিজে একবার চেষ্টা করো, তাতে সমাধানটা অনেক বেশি কাজে দেবে।
  \item \textbf{মন্তব্য} : পাশের কথা। কোনো সাবধানবাণী, কোনো সাধারণ ভুলের ইঙ্গিত,
        বা সামনের কোনো অধ্যায়ের দিকে ইশারা। বাদ দিলে হিসাব আটকাবে না, কিন্তু
        পড়লে বোঝাটা গভীর হবে।
  \item \textbf{সমস্যা} : তোমার কাজ।
\end{enumerate}
\end{idea}

\begin{idea}[{[N] এর অর্থ}] 
প্রতিটি সমস্যার বাঁ পাশে একটা সংখ্যা বসানো আছে। ওটা নম্বর নয়, কঠিনতার মাপ :

\begin{enumerate}
  \item $[1]$ সরাসরি প্রয়োগ। সংজ্ঞাটা জানলেই হয়ে যাবে।
  \item $[2]$ স্বাভাবিক। দুয়েকটা ধাপ জোড়া লাগাতে হবে।
  \item $[3]$ একটা বুদ্ধি লাগবে, যেটা সংজ্ঞা থেকে সরাসরি আসে না।
  \item $[4]$ কঠিন। সময় নিয়ে বসতে হবে।
  \item $[5]$ অলিম্পিয়াড পর্যায়ের।
\end{enumerate}

\noindent
অধ্যায়ের শুরুতে মোট পয়েন্ট বলা থাকে, যা সব চিহ্নের যোগফল।
\end{idea}

\begin{remark}
একটা পরামর্শ, যেটা এই সিরিজের বাকি সব পরামর্শের চেয়ে বেশি কাজে দেবে : কোনো
সমস্যায় আটকে গেলে সমাধান দেখার আগে অন্তত পনেরো মিনিট নিজে বসো। আটকে থাকা
অবস্থাটা সময় অপচয় নয়, ওটাই আসল শেখার জায়গা। সমাধান পড়ে ফেললে মনে হবে বুঝে গেছি,
কিন্তু পরেরবার একই সমস্যা আবার আটকাবে।

আর সমাধান দেখার পর খাতা বন্ধ করে নিজে আবার লিখে ফেলো। পড়তে পারা আর লিখতে পারা
এক জিনিস নয়।
\end{remark}

\section{সমস্যা}

\prob{1} নিচের কোনগুলো বিবৃতি (অর্থাৎ সত্য বা মিথ্যা হতে পারে), আর কোনগুলো নয়?
কারণ বলো।

\begin{enumerate}
  \item $17$ একটি মৌলিক সংখ্যা।
  \item $3x + 1$
  \item প্রতিটি জোড় সংখ্যা $2$ দ্বারা বিভাজ্য।
  \item ক্যালকুলাস পদার্থবিজ্ঞানের চেয়ে কঠিন।
  \item $x + 1 = 5$
\end{enumerate}

\begin{sol}
(ক) বিবৃতি, এবং সত্য।

(খ) বিবৃতি নয়। এটা একটা রাশি, কোনো দাবি করছে না। সত্য-মিথ্যা প্রশ্নটাই ওঠে না।

(গ) বিবৃতি, এবং সত্য।

(ঘ) বিবৃতি নয়। ``কঠিন'' কথাটার কোনো নির্দিষ্ট মানে দেওয়া হয়নি, তাই যাচাই করার
উপায় নেই। এটা মতামত।

(ঙ) $x$-এর মান না জানা পর্যন্ত এটা সত্যও নয় মিথ্যাও নয়। $x = 4$ হলে সত্য, নইলে
মিথ্যা। এ ধরনের বাক্যকে বলে খোলা বাক্য (open sentence); $x$-এ মান বসালে তবেই
বিবৃতি হয়।
\end{sol}

\prob{2} একজন শিক্ষার্থী দেখাতে চায় যে ``যেকোনো পূর্ণসংখ্যা $n$-এর জন্য
$n^2 - n$ জোড়''। সে লিখল :

\smallskip
\noindent
$n = 1$ হলে $n^2 - n = 0$, জোড়। $n = 2$ হলে $2$, জোড়। $n = 3$ হলে $6$, জোড়।
$n = 4$ হলে $12$, জোড়। $n = 5$ হলে $20$, জোড়। সুতরাং দাবিটি প্রমাণিত।

\smallskip
\noindent
এটা কি প্রমাণ? না হলে একটা সত্যিকারের প্রমাণ লেখো।

\begin{sol}
এটা প্রমাণ নয়। পাঁচটা মান পরীক্ষা করা হয়েছে, কিন্তু দাবিটা \emph{সব}
পূর্ণসংখ্যার কথা বলছে, আর পূর্ণসংখ্যা অসীম সংখ্যক। না-পরীক্ষা করা কোনো মানে
দাবিটা ভেঙে যেতে পারত।

\medskip
\textbf{প্রমাণ।} উৎপাদকে লিখি :
\[
  n^2 - n = n(n-1) .
\]
$n$ ও $n-1$ পরপর দুটি পূর্ণসংখ্যা। পরপর দুটি পূর্ণসংখ্যার একটি অবশ্যই জোড়,
কারণ জোড় আর বিজোড় পালা করে আসে। তাই গুণফলে একটি জোড় উৎপাদক আছে, অর্থাৎ
গুণফল $2$ দ্বারা বিভাজ্য।

লক্ষ করো, এখানে কোনো নির্দিষ্ট সংখ্যা বসাতে হয়নি। যুক্তিটা সব $n$-এর জন্য
একসাথে খেটেছে।
\end{sol}

\prob{2} নিচের প্রতিটি দাবি মিথ্যা। প্রতিটির জন্য একটি counterexample দাও।

\begin{enumerate}
  \item সব মৌলিক সংখ্যা বিজোড়।
  \item যেকোনো সংখ্যার বর্গ সেই সংখ্যার চেয়ে বড়।
  \item দুটি অমূলদ সংখ্যার যোগফল সবসময় অমূলদ।
\end{enumerate}

\begin{sol}
(ক) $2$। এটি মৌলিক এবং জোড়।

(খ) $n = 1$ নিলে $n^2 = 1$, বড় নয়, সমান। অথবা $n = \tfrac{1}{2}$ নিলে
$n^2 = \tfrac{1}{4}$, যা ছোট।

(গ) $\sqrt{2}$ ও $-\sqrt{2}$ দুটোই অমূলদ, কিন্তু যোগফল $0$, যা মূলদ।

তিনটে ক্ষেত্রেই একটামাত্র উদাহরণ দিয়েই কাজ শেষ। কোনো সাধারণ যুক্তির দরকার হয়নি।
\end{sol}

\prob{2} নিচের প্রতিটি ক্ষেত্রে $\Rightarrow$, $\Leftarrow$, নাকি
$\Leftrightarrow$ বসানো সৎ হবে বলো।

\begin{enumerate}
  \item $x = 5$ \qquad \framebox[1.4cm]{} \qquad $x^2 = 25$
  \item $n$ জোড় \qquad \framebox[1.4cm]{} \qquad $n^2$ জোড়
  \item $x > 3$ \qquad \framebox[1.4cm]{} \qquad $x > 1$
\end{enumerate}

\begin{sol}
(ক) $\Rightarrow$। $x = 5$ হলে অবশ্যই $x^2 = 25$। কিন্তু উল্টোটা নয়, কারণ
$x = -5$-ও $x^2 = 25$ দেয়।

(খ) $\Leftrightarrow$। $n$ জোড় হলে $n^2$ জোড়, এটা সরাসরি। আর $n$ বিজোড় হলে
$n^2$ বিজোড়, তাই $n^2$ জোড় হলে $n$ বিজোড় হতে পারে না, অর্থাৎ $n$ জোড়। দুই দিকই খাটে।

(গ) $\Rightarrow$। $3$-এর চেয়ে বড় হলে $1$-এর চেয়ে তো বড় বটেই। উল্টোটা নয়;
$x = 2$ ডান পাশ মানে কিন্তু বাঁ পাশ মানে না।
\end{sol}

\prob{2} নিচের ``সমাধান'' -এ কোন ধাপটা $\Leftrightarrow$ নয়, আর তার ফল কী হলো?

\smallskip
\noindent
সমাধান করো : $\sqrt{x + 6} = x$।
দুই পাশে বর্গ করি : $x + 6 = x^2$, অর্থাৎ $x^2 - x - 6 = 0$, অর্থাৎ
$(x-3)(x+2) = 0$। সুতরাং $x = 3$ অথবা $x = -2$।

\begin{sol}
দুই পাশে বর্গ করার ধাপটা $\Rightarrow$, $\Leftrightarrow$ নয়। কারণ $a = b$
হলে $a^2 = b^2$ ঠিকই, কিন্তু $a^2 = b^2$ হলে $a = b$ হতেই হবে না; $a = -b$-ও
হতে পারে।

ফলে যা পাওয়া গেল তা ``সম্ভাব্য উত্তর'', নিশ্চিত উত্তর নয়। যাচাই করি :

$x = 3$ : বাঁ পাশ $\sqrt{9} = 3$, ডান পাশ $3$। মিলেছে।

$x = -2$ : বাঁ পাশ $\sqrt{4} = 2$, ডান পাশ $-2$। মেলেনি। ($\sqrt{\;}$ চিহ্ন
অঋণাত্মক মান বোঝায়, তাই বাঁ পাশ কখনো ঋণাত্মক হতে পারে না।)

সুতরাং একমাত্র সমাধান $x = 3$। $x = -2$ বর্গ করার ধাপেই ঢুকে পড়েছিল, আসল
সমীকরণে ছিল না। এ ধরনের উত্তরকে বলে বহিরাগত মূল (extraneous root), আর এদের
ধরার একমাত্র উপায় যাচাই করা।
\end{sol}

\prob{1} কথায় লেখো :
(ক) $\forall n \in \mathbb{Z}, \; n^2 \geq 0$ \quad
(খ) $\exists x, \; x^2 = 2$।
এরপর বলো, দুটোর মধ্যে কোনটা একটামাত্র উদাহরণ দিয়ে প্রমাণ করা যায়।

\begin{sol}
(ক) ``প্রতিটি পূর্ণসংখ্যা $n$-এর জন্য $n^2$ অঋণাত্মক।''

(খ) ``এমন অন্তত একটি সংখ্যা $x$ আছে যার বর্গ $2$।''

দ্বিতীয়টা একটামাত্র উদাহরণ দিয়েই প্রমাণ করা যায় : $x = \sqrt{2}$ দেখিয়ে দিলেই
শেষ, কারণ দাবিটা শুধু ``অন্তত একটা আছে'' বলছে।

প্রথমটা উদাহরণ দিয়ে প্রমাণ করা যায় না, কারণ ``প্রতিটি'' মানে সব, আর সব
পূর্ণসংখ্যা গুনে শেষ করা যাবে না। এর জন্য সাধারণ যুক্তি লাগবে।

\medskip
মনে রাখার মতো নিয়ম : $\exists$-দাবি প্রমাণ করা সহজ, ভুল প্রমাণ করা কঠিন।
$\forall$-দাবির ঠিক উল্টো।
\end{sol}

\prob{2} দেখাও যে দুটি বিজোড় সংখ্যার যোগফল জোড়।

\begin{sol}
$n$ বিজোড় মানে $n = 2k + 1$, যেখানে $k$ একটি পূর্ণসংখ্যা।

ধরি $m$ ও $n$ দুটি বিজোড় সংখ্যা, অর্থাৎ $m = 2j + 1$ এবং $n = 2k + 1$ কোনো
পূর্ণসংখ্যা $j, k$-এর জন্য। তাহলে
\[
  m + n = (2j+1) + (2k+1) = 2j + 2k + 2 = 2(j + k + 1) .
\]
$j + k + 1$ একটি পূর্ণসংখ্যা, তাই $m+n$ জোড়।

\medskip
সাবধানবাণী : $m = 2j+1$ ও $n = 2j+1$ লেখা \emph{ভুল} হতো, কারণ ওতে ধরে নেওয়া
হয় সংখ্যা দুটি সমান। আলাদা সংখ্যার জন্য আলাদা অক্ষর দরকার।
\end{sol}

\prob{2} নিচের ``প্রমাণ'' -এর দোষ কী?

\smallskip
\noindent
দাবি : $\sqrt{2}$ অমূলদ।
প্রমাণ : যেহেতু $\sqrt{2}$-কে কোনো ভগ্নাংশ আকারে লেখা যায় না, তাই
$\sqrt{2}$ অমূলদ।

\begin{sol}
প্রমাণটা যা প্রমাণ করতে চাইছে, ঠিক সেটাকেই ধরে নিয়ে শুরু করেছে। ``ভগ্নাংশ
আকারে লেখা যায় না'' আর ``অমূলদ'' একই কথা, শুধু দুই ভাষায়। অর্থাৎ এখানে বলা
হয়েছে ``$\sqrt{2}$ অমূলদ, কারণ $\sqrt{2}$ অমূলদ''।

এ ধরনের যুক্তিকে বলে চক্রাকার যুক্তি (circular reasoning)। এতে কোনো নতুন
তথ্য যোগ হয় না, তাই এটা প্রমাণ নয়।

একটা সৎ প্রমাণকে দেখাতে হতো \emph{কেন} কোনো ভগ্নাংশ কাজ করে না, শুরু করতে হতো
অন্য কোনো মেনে নেওয়া সত্য থেকে। (সেই প্রমাণটা এই সিরিজে লাগবে না, তাই দেওয়া
হয়নি; ফলটা অধ্যায় ০-এ মেনে নেওয়া হয়েছে।)
\end{sol}

\end{document}
